\chapter{Модулярное соответствие в коммутанте и единый конечный след}

Предыдущая проверка восстановила канонический класс высоты, но оставила
неединственным унитарное отождествление семейного и вершинного пространств.
Физически необходимы не выбранные собственные векторы, а соответствие между
алгеброй спектральных проекторов модулярного оператора и алгеброй проекторов
вершин конечной геометрии.

\section{Алгебра вершин в полном коммутанте}

Пусть \(p_L,p_G,p_R\) являются ортогональными проекторами ранга три на
вершины частичной цепочки. Их полные вещественные пары обозначим
\begin{equation}
 \widehat p_i=p_i\oplus Jp_iJ^{-1}.
\end{equation}
На 18-мерном пространстве
\begin{equation}
 \widehat p_i^2=\widehat p_i,
 \qquad \widehat p_i\widehat p_j=0\quad(i\ne j),
 \qquad \widehat p_L+\widehat p_G+\widehat p_R=I_{18},
 \qquad \rank\widehat p_i=6.
 \label{eq:v5-commutant-vertex-projectors}
\end{equation}
Прямая проверка даёт
\begin{equation}
 [\widehat p_i,\pi(a)]=0,
 \qquad
 [\widehat p_i,J\pi(a)J^{-1}]=0
 \label{eq:v5-commutant-double-commutant}
\end{equation}
для полного базиса
\(\mathbb R\oplus M_3(\mathbb R)\oplus\mathbb C\).
Следовательно, проекторы порождают каноническую коммутативную алгебру
\begin{equation}
 \mathcal C_{\mathrm{vert}}
 =\operatorname{span}\{\widehat p_L,\widehat p_G,\widehat p_R\}
 \simeq\mathbb C^3
\end{equation}
в совместном коммутанте левого и противоположного действий.

\section{Соответствие спектральных проекторов}

Простой спектр старого факторного оператора порождает другую трёхточечную
алгебру
\begin{equation}
 \mathcal C_{\mathrm{mod}}=C^*(L_{\mathrm{fam}})\simeq\mathbb C^3.
\end{equation}
Её минимальная модулярная частота задаёт граф \(A_3\). Опора оператора
Дирака KO6 имеет тот же неориентированный граф. Сохраняющий смежность
\(*\)-изоморфизм
\begin{equation}
 \iota:\mathcal C_{\mathrm{mod}}
 \longrightarrow\mathcal C_{\mathrm{vert}}
 \label{eq:v5-commutant-iota}
\end{equation}
обязан перевести единственную вершину степени два в \(\widehat p_G\). Для концов
остаются только два варианта:
\begin{equation}
 (0,1,2)\longmapsto(L,G,R),
 \qquad
 (0,1,2)\longmapsto(R,G,L).
\end{equation}
Они индуцируют высоты
\begin{equation}
 (-1,0,+1),
 \qquad
 (+1,0,-1)
\end{equation}
и различаются общим обращением. Поэтому соответствие проекторных алгебр
единственно в физически существенном классе; независимые фазы собственных
векторов больше не входят в постановку.

\begin{theorem}[Каноническое проекторное соответствие]
Модулярный граф и граф опоры KO6 определяют единственный класс
\(*\)-изоморфизмов~\eqref{eq:v5-commutant-iota} с точностью до обращения
цепи. Этот класс не требует выбора базиса или унитарного сплетающего
оператора между двумя трёхмерными пространствами.
\end{theorem}

\section{Полная высота и верное состояние}

Выберем одного представителя на частичной половине и дополним посредством
вещественной структуры:
\begin{equation}
 h_p=-p_L+p_R,
 \qquad
 h_F=h_p\oplus(-h_p).
\end{equation}
Тогда
\begin{equation}
 Jh_FJ^{-1}=-h_F,
 \qquad
 [h_F,\Gamma_F]=0,
 \qquad
 [h_F,\pi(a)]=[h_F,J\pi(a)J^{-1}]=0.
\end{equation}
Все машинные дефекты равны нулю.

На полной матричной алгебре
\begin{equation}
 \mathcal B_{\mathrm{par}}=M_{18}(\mathbb C)
\end{equation}
определим для любого \(\beta>0\) состояние
\begin{equation}
 \rho_\beta
 =\frac{e^{-\beta h_F}}{\Tr e^{-\beta h_F}}.
 \label{eq:v5-commutant-rho}
\end{equation}
Оно положительно, верно и коммутирует с обоими действиями физической
алгебры. Параметр \(\beta\) меняет абсолютные вероятности уровней, но не
знак модулярных частот и не выбранную ориентацию. Поэтому никакое
наблюдаемое число по \(\beta\) здесь не вычисляется.

\section{Модулярная поляризация оператора Дирака}

Для состояния~\eqref{eq:v5-commutant-rho}
\begin{equation}
 \Delta_\rho(A)=\rho_\beta A\rho_\beta^{-1},
 \qquad
 K_{\text{мод}}=-\log\Delta_\rho
 =\beta\operatorname{ad}_{h_F}.
\end{equation}
Пусть \(\Pi_{+1}\) обозначает спектральный проектор
\(\operatorname{ad}_{h_F}\) на собственное значение \(+1\). Тогда
ориентированный дифференциал определяется уже не вручную, а формулой
\begin{equation}
 d:=\Pi_{+1}(D_F),
 \qquad
 D_F=d+d^\dagger.
 \label{eq:v5-commutant-modular-d}
\end{equation}
Вещественная структура обращает частоту:
\begin{equation}
 JdJ^{-1}=d^\dagger.
\end{equation}
На полном явном представлении получены нулевые дефекты самосопряжённости,
нечётности по \(\Gamma_F\), вещественности, условия нулевого порядка и
условия первого порядка. Дефект равенства
\(K_{\text{мод}}=\beta\operatorname{ad}_{h_F}\) на всех ненулевых блоках
\(d\) не превышает \(1.2\cdot10^{-16}\).

\begin{proposition}[Модулярное происхождение ориентации]
Положительно-частотная часть полного оператора Дирака совпадает с
повышающим дифференциалом согласованной цепи. Ориентация определяется
состоянием до задания матриц \(X,Y\) и не является знаком, вставленным в
квартичный функционал.
\end{proposition}

\section{Кривизна и единственный след}

На средней вершине имеем точное тождество
\begin{equation}
 [d,d^\dagger]\big|_{V_G}
 =XX^\dagger-Y^\dagger Y.
\end{equation}
Полная матричная алгебра имеет единственный нормированный след
\begin{equation}
 \tau_{18}(A)=\frac1{18}\Tr_{18}A.
\end{equation}
После условной нормировки на шестимерном среднем секторе
\begin{equation}
 \frac16\Tr_{18}\!\left(
 \widehat p_G[d,d^\dagger]^2
 \right)
 =\tau_3(XX^\dagger-Y^\dagger Y)^2.
 \label{eq:v5-commutant-trace-identity}
\end{equation}
Численный дефект~\eqref{eq:v5-commutant-trace-identity} равен
\(3.6\cdot10^{-15}\).

\section{Почему прежний запрет скрученного следа не применяется}

Предыдущий запрет относился к автоморфизму, переставляющему центральные
идемпотенты разных рангов. Здесь центр \(M_{18}(\mathbb C)\) одномерен, а
проекторы \(\widehat p_L,\widehat p_G,\widehat p_R\) не центральны в полной матричной алгебре. Они
лишь принадлежат коммутанту физического представления. Состояние
\(\rho_\beta\) не переставляет центральные блоки: оно поляризует
нецентральные переходные операторы. Явный коммутатор с межвершинной
матричной единицей имеет норму \(\sqrt3\ne0\).

Следовательно, новая конструкция не опровергает теорему о невозможности
верного следа для центрального обмена и не пользуется запрещённым
скручиванием.

\section{Граница результата}

Получен первый положительный родительский подобъект Тома V:
\begin{equation}
 \boxed{
 M_{18}(\mathbb C),\ \rho_\beta,\ \tau_{18},\ J,\ \Gamma_F,\ D_F
 }
\end{equation}
в одной конструкции дают верное состояние, ориентацию, KO6 и точный
квадрат отображения момента.

Но это ещё не полная родительская архитектура наблюдаемого сектора. В
данном объекте пока отсутствуют:
\begin{enumerate}
 \item представление Стандартной модели и существующий единый след
 \(M_{60}(\mathbb C)\) без повторного счёта поколенческого триплета;
 \item одно полное действие для калибровочного и семейного секторов;
 \item физический фактор БВ/БРСТ и полный гессиан;
 \item две независимые слепые наблюдаемые.
\end{enumerate}

\begin{nogocriterion}[Запрет формального тензорного раздувания]
Нельзя объявить замыкание простым тензорным произведением \(M_{18}\) и
\(M_{60}\). Следующая конструкция обязана показать, какие три степени
свободы являются поколениями, какие являются вершинами, и исключить их
двойной счёт. Любой новый центральный вес или независимый коэффициент
между калибровочным и семейным действиями закрывает путь.
\end{nogocriterion}

\section{Вердикт}

\begin{protocol}
\begin{enumerate}
 \item проекторное модулярно-вершинное соответствие:
 \textbf{получено};
 \item один полный конечный носитель для состояния и KO6:
 \textbf{получен};
 \item верное положительное состояние:
 \textbf{получено};
 \item ориентация до задания \(X,Y\):
 \textbf{получена};
 \item единственный полный матричный след:
 \textbf{получен};
 \item точный квадрат отображения момента:
 \textbf{получен};
 \item объединение со Стандартной моделью и следом \(M_{60}\):
 \textbf{не выполнено};
 \item полное действие и физический гессиан:
 \textbf{не получены};
 \item локальная родительская подобархитектура:
 \textbf{прошла проверку};
 \item физическое замыкание:
 \textbf{не достигнуто}.
\end{enumerate}
\end{protocol}

Следующая проверка должна объединить эту конструкцию с ранее полученным
семейно-калибровочным следом \(M_{60}\), не удваивая поколенческий триплет и
не вводя нового относительного веса.

Машинная проверка и сертификат сохранены в
\path{s2t_v5_modular_commutant_parent_correspondence_gate.py} и
\path{s2t_v5_modular_commutant_parent_correspondence_gate_results.json}.